\author{Lothar Schmid}
\documentclass[12pt]{../../../inc/mybib}

\setincpath{../../../inc/}

\usepackage{bible_style}
\usepackage{bibeltext}
\usepackage{markierungen}
\graphicspath{{../../../assets/images/}}
\usepackage{header}


% ensure scrlayer-scrpage has sufficient footheight
\setlength{\footheight}{20.4pt}
\newcommand{\doublehline}{\hline\noalign{\vskip 0.8mm}\hline}
\begin{document}

\section{Tabellarische Auflisung der Konjunktionen im Nehemiabrief}
\renewcommand{\arraystretch}{1.8}
\begin{longtable}{|p{1cm} p{3cm} p{2cm} p{8cm}|}
\hline 
\rowcolor{gray!20}

\textbf{Vers} & \textbf{Konjunktion} & \textbf{Art} & \textbf{Bemerkung} \\ 

\hline
\endfirsthead

\hline 
\rowcolor{gray!20}
\textbf{Vers} & \textbf{Konjunktion} & \textbf{Art} & \textbf{Bemerkung} \\ 
\hline
\endhead
\verslink{Nehemia}{1}{2}
& Da kam \dots
& \textit{Grund}
& Als Nehemia auf der Burg Susan war.
\\ 
\hline
\verslink{Nehemia}{1}{6}
& Sowohl ich als auch \dots
& \textit{Vergleich}
& Nehemia gibt die Schuld nicht nur seinen Vorfahren, sondern fühlt sich mitschuldig.
\\ 
\hline
\verslink{Nehemia}{1}{8}
& Ihr, wenn \dots
& \textit{Bedingung}
& Das ganze Volk muss umkehren, dann wird Gott es wieder zurückbringen.
\\ 
\hline
\verslink{Nehemia}{1}{11}
& Ach, \dots
& \textit{Ausruf}
& Verzweiflungsausdruck und Bitte an Gott ihm zu helfen.
\\ 
\doublehline
\verslink{Nehemia}{2}{1}
& Doch, \dots
& \textit{Ergebnis}
& Der König hatte erkannt, das mit Nehemia etwas nicht stimmt. Nehemia hat Gott um eine Gelegenheit gebeten, um mit dem König sprechen zu können und diese Gelegenheit hat er nun bekommen. 
\\ 
\hline
\verslink{Nehemia}{2}{7}
& wenn \dots
& \textit{Art}
& Anfrage an den König. Seine Bestätigung.
\\ 
\hline
\verslink{Nehemia}{2}{7}
& dass \dots
& \textit{Zweck}
& Für was der Brief helfen soll. Sichere Durchreise und Baumaterial.
\\ 
\hline
\verslink{Nehemia}{2}{8}
& dass \dots
& \textit{Zweck}
& Der Grund des Briefes
\\ 
\hline
\verslink{Nehemia}{2}{10}
& dass \dots
& \textit{Folge}
& Verwunderung wieso Nehemia sich dem Volk annimmt.
\\ 
\hline
\verslink{Nehemia}{2}{11}
& Als\dots
& \textit{Zeit}
& Nehemia kommt in Jerusalem an und geht direkt an die Arbeit. Er will sich ein Bild von der Lage machen, bevor er mit anderen darüber spricht.
\\ 
\hline
\verslink{Nehemia}{2}{16}
& denn
& \textit{Grund}
& Nehemia hat mit den Ansässigen noch nicht über sein vorhaben gesprochen.
\\ 
\hline
\verslink{Nehemia}{2}{17}
& Da
& \textit{Zeit}
& Nehemia gibt seine Beobachtungen weiter und ruft zum Handeln auf.
\\ 
\hline
\verslink{Nehemia}{2}{18}
& dass
& \textit{Folge}
& Die Hand Gottes und der König waren ihm wohlgesinnt
\\ 
\hline
\verslink{Nehemia}{2}{18}
& da
& \textit{Folge}
& Folge aus der Zuversicht von Nehemia
\\ 
\hline
\verslink{Nehemia}{2}{20}
& Doch\dots
& \textit{Ein\-schränkung}
& Nicht sie sind die Besitzer von Jerusalem sondern Gott und dieser ist es der hilft.
\\ 
\doublehline
\verslink{Nehemia}{3}{33}
& dass
& \textit{Grund}
& Der Feind hat gehört was die Einwohner in Israel machen.
\\ 
\hline
\verslink{Nehemia}{3}{37}
& denn
& \textit{Grund}
& Weil der Feind das Volk demotiviert hat, soll Gott eingreifen.
\\ 
\doublehline
\verslink{Nehemia}{4}{1}
& 2x dass
& \textit{Grund}
& Der Feind hat hat wieder mitbekommen, dass sie weiter an der Mauer bauen.
\\ 
\hline
\verslink{Nehemia}{4}{2}
& um
& \textit{Grund}
& Sie rotten sich zusammen um gegen das Volk zu kämpfen
\\ 
\hline
\verslink{Nehemia}{4}{3}
& Da
& \textit{Grund}
& Nehemia und das Volk holt unterstützung  von Gott.
\\ 
\hline
\verslink{Nehemia}{4}{5}
& bis
& \textit{Zeit}
& Der Angriff soll verheimlicht sein, bis sie mitten und den Bauenden sind
\\ 
\hline
\verslink{Nehemia}{4}{7}
& Da
& \textit{Grund}
& Nehemia hat mitbekommen, dass sie Angegriffen werden sollen.
\\ 
\hline
\verslink{Nehemia}{4}{9}
& 2x dass
& \textit{Grund}
& Die Feinde haben mitbekommen, dass Nehemia mitbekommen hat, dass sie angreifen wollen.
\\ 
\hline
\verslink{Nehemia}{4}{9}
& da
& \textit{Grund}
& Darum kehrten sie alle zum Schutz der Mauer zurück zur Arbeit.
\\ 
\hline
\verslink{Nehemia}{4}{11}
& 2x mit
& \textit{Grund}
& Bei der Arbeit trugen sie zusätzlich zur Verteidigung noch eine Waffe bei sich
\\ 
\hline
\verslink{Nehemia}{4}{15}
& Und so
& \textit{Ergebnis}
& Mit der Zuversicht, dass Gott Kämpfen wird und mit der Waffe in der Hand
\\ 
\hline
\verslink{Nehemia}{4}{16}
& Auch
& \textit{Fort\-setzung}
& Nehemia will dass auch in der Nacht die Sicherheit gewährleistet ist.
\\ 
\hline
\verslink{Nehemia}{4}{16}
& jederman
& \textit{Zahl}
& Nehemia will, dass auch in der Nacht die Sicherheit gewährleistet ist.
\\ 
\hline
\verslink{Nehemia}{4}{16}
& so dass
& \textit{Ergebnis}
& Tag und Nacht soll auf die Sicherheit geachtet werden.
\\ 
\doublehline
\verslink{Nehemia}{5}{2}
& dass
& \textit{Zweck}
& Weil manche viele Kinder hatten, möchten sie jetzt mehr Getreide bekommen. Weil diese auf dem Bau sind.
\\ 
\hline
\verslink{Nehemia}{5}{3}
& damit
& \textit{Zweck}
& Manche mussten ihre Felder verpfänden und möchten jetzt mehr Getreide
\\ 
\hline
\verslink{Nehemia}{5}{5}
& und jetzt
& \textit{Zeit}
& Obwohl sie doch vom gleichen Volk sind, müssen ihre Söhne und Töchter Sklavenarbeit verrichten.
\\ 
\hline
\verslink{Nehemia}{5}{6}
& und
& \textit{Fort\-setzung}
& Nehemia wurde wegen den Aussagen der Menschen zornig
\\ 
\hline
\verslink{Nehemia}{5}{8}
& welche
& \textit{Ursache}
& Erklärung von Nehemia, dass sie die Rückführung aller Brüder bewerkstelligt haben
\\ 
\hline
\verslink{Nehemia}{5}{8}
& soweit
& \textit{Mittel}
& Soweit sie die Möglichkeit hatten
\\ 
\hline
\verslink{Nehemia}{5}{9}
& solltet
& \textit{Grund}
& Das Volk soll zusammenhalten
\\ 
\hline
\verslink{Nehemia}{5}{10}
& Aber
& \textit{Ergebnis}
& Auch Nehemia hat geholfen mit Geld und Getreide aus seinem eigenen Sack.
\\ 
\hline
\verslink{Nehemia}{5}{13}
& So
& \textit{Folge}
& Wenn sich jemand weigern sollte zu helfen, soll Gott ihn richten.
\\ 
\hline
\verslink{Nehemia}{5}{14}
& Auch
& \textit{Folge}
& Nehemia hat sich auch nachdem er Stadthalter wurde immer selbst versorgt.
\\ 
\hline
\verslink{Nehemia}{5}{15}
& Aber
& \textit{Gegen\-satz}
& Die früheren Stadthalter haben das Volk ausgenutzt
\\ 
\hline
\verslink{Nehemia}{5}{15}
& Auch
& \textit{Erklärung}
& Sogar die Diener der Stadthalter haben mitgemacht
\\ 
\hline
\verslink{Nehemia}{5}{15}
& Und trotzdem
& \textit{Erklärung}
& Obwohl er alles aus seinem eignen Sack bezahlt fordert er keinen Lohn für sich.
\\ 
\doublehline
\verslink{Nehemia}{6}{1}
& dass
& \textit{Grund}
& Die Feinde hörten, dass die Mauer bald fertig ist.
\\ 
\hline
\verslink{Nehemia}{6}{2}
& da
& \textit{Grund}
& Weil die Mauer trotz der kriegerischen Bedrohung fertig wurde, wird nun List und Verrat angewandt.
\\ 
\hline
\verslink{Nehemia}{6}{3}
& Warum
& \textit{Grund}
& Frage nach Grund wieso er das Werk stehen lassen. Nehemia hat den Verrat geahnt
\\ 
\hline
\verslink{Nehemia}{6}{6}
& dass
& \textit{Ziel}
& Als Grund für den Mauerbau wird Nehemia beschuldigt einen Aufstand zu planen
\\ 
\hline
\verslink{Nehemia}{6}{6}
& darum
& \textit{Grund}
& Der Aufstand und um König zu werden ist der Grund des Mauerbau
\\ 
\hline
\verslink{Nehemia}{6}{7}
& Auch \dots damit
& \textit{Zweck}
& Nehemia soll Propheten bestellt haben, um sich selber zum König zu machen.
\\ 
\hline
\verslink{Nehemia}{6}{8}
& sondern
& \textit{Gegen\-satz}
& Nehemia macht klar, dass er keine solche Absichten hat und alles erfunden ist.
\\ 
\hline
\verslink{Nehemia}{6}{9}
& denn
& \textit{Grund}
& Der Brief ist nur da zum Einschüchtern
\\ 
\hline
\verslink{Nehemia}{6}{10}
& denn
& \textit{Grund}
& Nehemia wird mit einem Vorhand gelockt um ihn zu töten
\\ 
\hline
\verslink{Nehemia}{6}{10}
& Aber
& \textit{Folge}
& Nehemia erkennnt der Verrat
\\ 
\hline
\verslink{Nehemia}{6}{13}
& Dazu \dots damit
& \textit{Grund}
& Nehemia erkennnt den Grund wieso dieser Mann zu ihm kam.
\\ 
\hline
\verslink{Nehemia}{6}{16}
& Da
& \textit{Folge}
& Die Völker erkennen, dass da Gott mit im Spiel sein muss.
\\ 
\hline
\verslink{Nehemia}{6}{17}
& Auch
& \textit{Mittel}
& Viele liefen zum Feind über
\\ 
\hline
\verslink{Nehemia}{6}{18}
& denn
& \textit{Ursache}
& Sie waren mit dem Feind eidlich verbunden.
\\ 
\hline
\verslink{Nehemia}{6}{19}
& Auch
& \textit{Grund}
& Sie pflegten von seinen guten Taten zu reden.
\\ 
\doublehline
\verslink{Nehemia}{7}{1}
& da
& \textit{Zeit}
& Es werden die Türflügel eingesetzt
\\ 
\hline
\verslink{Nehemia}{7}{2}
& denn
& \textit{Grund}
& Chanani wurde eingesetzt weil er ein treuer Mann war.
\\ 
\hline
\verslink{Nehemia}{7}{5}
& da
& \textit{Ursache}
& Gott gab Nehemia den Auftrag die Edlen zu versammeln
\\ 
\hline
\verslink{Nehemia}{7}{5}
& um
& \textit{Grund}
& Um sie in das Geschlechtsregister einzuschreiben
\\ 
\hline
\verslink{Nehemia}{7}{72}
& und
& \textit{Zweck}
& Der siebte Monat, der Neujahrstag
\\ 
\doublehline
\verslink{Nehemia}{8}{1}
& da
& \textit{Grund}
& An diesem Monat versammelte sich das ganze Volk
\\ 
\hline
\verslink{Nehemia}{8}{1}
& dass
& \textit{Grund}
& Esra soll das Gebetsbuch mitbringen
\\ 
\hline
\verslink{Nehemia}{8}{1}
& welches
& \textit{Hin\-weis}
& Es ist die Tora. Die Gesetze von Gott die Gott Mose gegeben hat.
\\ 
\hline
\verslink{Nehemia}{8}{2}
& Da
& \textit{Grund}
& Weil das Volk ihn bat, bringt Esra das Gesetz.
\\ 
\hline
\verslink{Nehemia}{8}{5}
& Denn
& \textit{Grund}
& Esra stand ehöht zum Volk darum sah das ganze Volk, dass er das Buch öffnet.
\\ 
\hline
\verslink{Nehemia}{8}{10}
& Denn
& \textit{Grund}
& Sie sollen alle am Fest beteiligen, weil dieser Tag vor dem Herrn Heilig. Auch Arme und Alte sollen mitfeiern können.
\\ 
\hline
\verslink{Nehemia}{8}{10}
& Denn
& \textit{Grund}
& Esra sieht, dass das Volk ehrlich eine Freude an Gott hat.
\\ 
\hline
\verslink{Nehemia}{8}{11}
& Denn
& \textit{Grund}
& Es soll keine Gelage werden, das Fest soll ehrfürchtig begangen werden weil dieser Tag heilig ist.
\\ 
\hline
\verslink{Nehemia}{8}{12}
& Da
& \textit{Grund}
& Aus dem Grund weil sie Gott lieben und Freude haben, ging das ganze Volk hin, verschenkte Essen und machten ein grosses Fest.
\\ 
\hline
\verslink{Nehemia}{8}{12}
& Denn
& \textit{Grund}
& Sie hatten die Worte von Esra verstanden, darum machten sie das Fest und gaben auch den Armen und Schwachen.
\\ 
\hline
\verslink{Nehemia}{8}{13}
& Um
& \textit{Grund}
& Am zweiten Tag versammelten sich die Elite und studierten selber nochmals das Wort.
\\ 
\hline
\verslink{Nehemia}{8}{14}
& dass
& \textit{Zweck}
& Sie fanden in den Schriften das Laubhüttenfest. Dieses Fest war genau in der Mitte dieses Monats angelegt.
\\ 
\hline
\verslink{Nehemia}{8}{15}
& und dass
& \textit{Fort\-setzung}
& Sie verkündigten über ganz Isreal, dass das Volk dieses Laubhüttenfest halten soll
\\ 
\hline
\verslink{Nehemia}{8}{17}
& denn
& \textit{Grund}
& Das Laubhüttenfest wurde seit Josua nicht mehr durchgeführt. Es wurde darum auch zu einer grossen Freude.
\\ 
\doublehline
\verslink{Nehemia}{9}{2}
& Und
& \textit{Folge}
& Nach diesen Feiern sonderte sich Isreal von den Fremden ab und bekannten ihre Sünden.
\\ 
\hline
\verslink{Nehemia}{9}{5}
& Und dann
& \textit{Ergebnis}
& Das Ergebnis war, dass die Leviten anfingen Gott zu preisen und loben. Sie machten ein Dankgebet an den Herrn.
\\ 
\hline
\verslink{Nehemia}{9}{32}
& Und nun
& \textit{Folge}
& Nach dem Gebet wollen die Leviten ihren Bund mit Gott erneuern. Zuerst machen sie Gott auf den Bund aufmerksam und merken an, dass sie diesen gebrochen haben.
\\ 
\hline
\verslink{Nehemia}{9}{33}
& Doch
& \textit{Mittel}
& Sie weisen Gott auf seine Gerechtigkeit hin. Eine Gerechtigkeit die auch schon in ganzen Gebet, als sie die Geschichte Israel vortrugen, zum Vorschein kam.
\\ 
\hline
\verslink{Nehemia}{9}{33}
& Denn
& \textit{Grund}
& Sie sehen ihren Fehler ein.
\\ 
\doublehline
\verslink{Nehemia}{10}{1}
& Und bei
& \textit{Ergebnis}
& Sie schrieben eine Vertrag, der von der Elite unterschrieben werden soll. Mit diesem Vertrag wollen sie wieder nach dem Gesetz des Herrn leben.
\\ 
\doublehline
\verslink{Nehemia}{11}{1}
& Und
& \textit{Fort\-setzung}
& Nach dem aufsetzten und unterschreiben des Vertrags teilten sie die Arbeiten unter dem Volk auf.
\\ 
\doublehline
\verslink{Nehemia}{12}{27}
& Und \dots um
& \textit{Grund}
& Alle Leviten aus Israel wurden zur Einweihung der Mauer eingeladen
\\ 
\hline
\verslink{Nehemia}{12}{28}
& Sowohl
& \textit{Vergleich}
& Nicht nur Sänger aus Jerusalem sondern aus ganz Israel haben sich dort versammelt.
\\ 
\hline
\verslink{Nehemia}{12}{29}
& denn
& \textit{Folge}
& Die Sänger campierten draussen vor der Mauer.
\\ 
\hline
\verslink{Nehemia}{12}{43}
& und auch
& \textit{Folge}
& Auch die Frauen und Kinder freuten sich. Die Freude wurde weithin vernommen.
\\ 
\hline
\verslink{Nehemia}{12}{44}
& dass
& \textit{Ziel}
& Es wird mit der Umsetzung der Gesetze angefangen
\\ 
\hline
\verslink{Nehemia}{12}{44}
& denn
& \textit{Grund}
& Juda, also die Juden um Jerusalem freuten sich an dem Dienst der Leviten.
\\ 
\hline
\verslink{Nehemia}{12}{44}
& denn \dots seit
& \textit{Grund}
& Schon seit der Zeit Davids gab es schon Sänger und Chöre.
\\ 
\doublehline
\verslink{Nehemia}{13}{1}
& An jenem
& \textit{Zeit}
& Es wurde jetzt regemässig in den Gesezten gelesen 
\\ 
\hline
\verslink{Nehemia}{13}{1}
& dass
& \textit{Folge}
& und fanden heraus, dass die Ammoniter und Moabiter nicht in die Versammlungen mit einbeziehen dürfen.
\\ 
\hline
\verslink{Nehemia}{13}{2}
& denn
& \textit{Grund}
& Durch den Verrat von Bileam.
\\ 
\hline
\verslink{Nehemia}{13}{3}
& da
& \textit{Folge}
& Israel trennte sich von dem Mischvolk.
\\ 
\hline
\verslink{Nehemia}{13}{7}
& indem
& \textit{Mittel}
& Der Priester hat sich ein eigenes kleines Reich geschaffen, in dem er seine Gaunerei betreiben konnte
\\ 
\verslink{Nehemia}{13}{8}
& und es missfiel
& \textit{Folge}
& Nehemia liess die Kammer wieder räumen und reinigen, so dass so wieder als Speisekammer für die Opfer gebraucht werden konnte.
\\ 
\hline
\verslink{Nehemia}{13}{11}
& Da
& \textit{Grund}
& Es war nicht nur der Priester in die Korruption involwiert sondern auch noch mehrer Leviten. Darum hat Nehemia die Zustände gegenüber dem Vorsteher angesprochen
\\ 
\hline
\verslink{Nehemia}{13}{13}
& Denn
& \textit{Grund}
& Nehemia hat die Priester und Leviten als treu erachtet und vertraute auf sie.
\\ 
\hline
\verslink{Nehemia}{13}{15}
& Denn
& \textit{Zeit}
& Zur gleichen Zeit erkannte er nicht nur dass Korruption unter den Leviten herrschte, sondern dass auch das Volk von den Gesetzen die sie gelesen hatten abwichen. Sie hielten den Sabbat nicht mehr ein.
\\ 
\hline
\verslink{Nehemia}{13}{16}
& Auch
& \textit{Bestätigung}
& Auch Ausländer machten am Sabbat ihre Arbeit Obwohl alle Einwohner sich an die Sabbatgebote halten müssen.
\\ 
\hline
\verslink{Nehemia}{13}{19}
& dass \dots bis
& \textit{Folge}
& Nehemia gibt Anweisung am Sabbat die Tore geschlossen zu halten
\\ 
\hline
\verslink{Nehemia}{13}{19}
& damit
& \textit{Grund}
& Das keine Waren und Lasten mehr transportiert werden.
\\ 
\hline
\verslink{Nehemia}{13}{21}
& Wenn
& \textit{Drohung}
& Nehemia war sogar bereit Gewalt anzuwenden, wenn sie sich nicht an die Gebote halten.
\\ 
\hline
\verslink{Nehemia}{13}{22}
& dass
& \textit{Art}
& Sie sollen sich am Sabbat waschen, heiligen und das Tor bewachen und den Sabbat heiligen.
\\ 
\hline
\verslink{Nehemia}{13}{22}
& Auch das
& \textit{Hinweis}
& Nehemia bittet Gott, dass er die Verfehlungen ihm nicht anrechnet.
\\ 
\hline
\verslink{Nehemia}{13}{23}
& Auch sah
& \textit{Fort\-setzung}
& Nehemia sieht, dass wieder Fremdlinge Einzug hielten. 
\\ 
\hline
\verslink{Nehemia}{13}{25}
& Wenn
& \textit{Folge}
& Er erklärt was für Folgen dieses verhalten haben kann. Er verweist auf die Geschichte des Volkes.
\\ 
\hline
\verslink{Nehemia}{13}{30}
& Und
& \textit{Folge}
& Er stellte wieder die Ordnung her, schmiss die Fremdlinge raus, setzte die Dienstordnung für Priester und Leviten wieder ein. 
\\ 
\hline

\end{longtable}
Gedenke es mir, mein Gott, diese Liste zum Guten.
\end{document}